\chapter{Multiarmed Bandit Problem (MAB)}
\section{MAB}

\subsection{Il Multi-Armed Bandit Problem e il dilemma Exploration--Exploitation}

Il materiale introduce il \emph{Multi-Armed Bandit Problem} (MAB), con particolare attenzione alle applicazioni Web e ai sistemi di raccomandazione (\emph{Recommender Systems}, abbreviati nel seguito come RecSys).

Il problema fondamentale alla base dei MAB è il cosiddetto \emph{exploration--exploitation dilemma}, cioè il problema di stabilire quanto un agente debba:
\begin{itemize}
    \item sfruttare la conoscenza già acquisita, scegliendo l'opzione che al momento sembra migliore (\emph{exploitation});
    \item esplorare opzioni differenti, anche apparentemente peggiori, allo scopo di acquisire nuova informazione (\emph{exploration}).
\end{itemize}

Il dilemma nasce perché le conseguenze delle diverse azioni non sono note a priori con certezza. L'agente deve quindi contemporaneamente \emph{imparare} quali azioni siano migliori e \emph{decidere} quali azioni eseguire.

\subsubsection{Esempio motivazionale nei sistemi di raccomandazione}

Si consideri un utente che utilizza per la prima volta un'applicazione di food delivery. L'utente seleziona un ristorante messicano e assegna alla propria esperienza una valutazione molto positiva.

Un sistema di raccomandazione può dedurre da questa prima interazione che l'utente gradisca la cucina messicana. Di conseguenza, comincia a raccomandare altri ristoranti messicani. Se l'utente prova un secondo ristorante messicano e assegna nuovamente una valutazione positiva, il sistema rafforza ulteriormente la propria convinzione.


Si instaura quindi il seguente meccanismo:
\begin{enumerate}
    \item l'utente manifesta inizialmente una preferenza;
    \item il RecSys utilizza tale preferenza per formulare raccomandazioni simili;
    \item tali raccomandazioni generano ulteriori osservazioni dello stesso tipo;
    \item i nuovi feedback confermano il modello già costruito;
    \item il sistema diventa progressivamente più orientato verso la stessa tipologia di contenuto.
\end{enumerate}

La strategia funziona inizialmente, ma può diventare problematica. L'utente può infatti stancarsi della cucina messicana e, ignorando le raccomandazioni, provare autonomamente un ristorante italiano, scoprendo di apprezzarlo moltissimo.

Il problema è che il RecSys dispone ormai di molti feedback relativi alla cucina messicana. Una singola osservazione positiva sulla cucina italiana potrebbe non essere sufficiente a modificare rapidamente il modello delle preferenze costruito fino a quel momento.

Questo esempio mostra il limite di una strategia eccessivamente orientata all'\emph{exploitation}: sfruttando continuamente ciò che sembra già funzionare, il sistema raccoglie poche informazioni sulle alternative e può quindi restare intrappolato in una rappresentazione incompleta delle preferenze dell'utente.

\subsubsection{Esempio motivazionale: scelta di un banner pubblicitario}

Un secondo esempio riguarda una campagna pubblicitaria online nella quale si dispone di differenti versioni di un banner.

L'obiettivo è scegliere il banner più efficace, cioè quello in grado di massimizzare il \emph{Click-Through Rate} (CTR), ossia la frequenza con cui il banner visualizzato viene cliccato.

Per individuare il banner migliore è possibile effettuare una fase sperimentale:
\begin{itemize}
    \item a un gruppo di utenti viene mostrato il primo banner;
    \item ad altri utenti viene mostrato il secondo;
    \item il procedimento viene ripetuto per tutte le versioni disponibili;
    \item vengono infine confrontati i rispettivi valori di conversione, nel caso considerato il CTR.
\end{itemize}

La fase esplorativa è necessaria per ottenere informazione sufficiente a distinguere le diverse alternative. Tuttavia, ogni volta che viene mostrato un banner sub-ottimale, si incorre potenzialmente in una perdita di profitto rispetto a ciò che si sarebbe ottenuto mostrando il banner migliore.

Esiste dunque una tensione fondamentale:
\begin{itemize}
    \item esplorare troppo poco può impedire di identificare il banner realmente migliore;
    \item esplorare troppo a lungo, quando si dispone già di sufficiente evidenza sull'alternativa migliore, diventa controproducente.
\end{itemize}

\subsubsection{Definizione di exploitation ed exploration}

Nel contesto delle slide, i due concetti fondamentali sono definiti come segue.

\paragraph{Exploitation.}
Consiste nello scegliere l'opzione ritenuta migliore sulla base delle conoscenze acquisite fino al momento corrente.

Una politica fortemente \emph{exploitive} tende quindi a privilegiare il guadagno a breve termine. La forma estrema di questa idea è un approccio \emph{greedy}: in ogni istante viene scelta semplicemente l'opzione che, sulla base dei dati disponibili, sembra migliore.

Il rischio è che una scelta iniziale errata venga rinforzata nel tempo e conduca a una soluzione sub-ottimale nel lungo periodo.

\paragraph{Exploration.}
Consiste nello scegliere deliberatamente un'opzione diversa da quella attualmente ritenuta migliore, allo scopo di acquisire nuova conoscenza.

Una politica fortemente \emph{explorative} può quindi sacrificare reward nel breve termine per aumentare la probabilità di individuare una soluzione migliore nel lungo periodo.

\paragraph{Obiettivo del MAB.}
Il Multi-Armed Bandit Problem cerca di definire politiche \emph{principled}, cioè fondate su criteri sistematici, in grado di trovare un opportuno \emph{trade-off} tra exploration ed exploitation.

\subsubsection{Interpretazione generale del dilemma}

Il dilemma exploration--exploitation si manifesta in tutti i contesti nei quali:
\begin{enumerate}
    \item un agente possiede più azioni o decisioni possibili;
    \item l'esito delle azioni è incerto;
    \item l'agente può utilizzare gli esiti osservati per migliorare le decisioni successive.
\end{enumerate}

Il problema combina pertanto due componenti:
\begin{itemize}
    \item \emph{decision making}, perché l'agente deve prendere decisioni;
    \item \emph{learning}, perché la qualità delle decisioni dipende dall'informazione raccolta durante l'interazione.
\end{itemize}

Le slide richiamano anche un'analogia con il concetto psicologico di \emph{chiusura cognitiva}:
\begin{itemize}
    \item un soggetto eccessivamente ``chiuso'' tende a esplorare poco e ad affidarsi rapidamente alla prima soluzione che appare soddisfacente;
    \item un soggetto eccessivamente ``aperto'' tende invece a continuare a esplorare senza convergere verso una soluzione definitiva.
\end{itemize}

Anche in questo caso il comportamento efficace richiede un equilibrio.

\subsection{Exploration ed exploitation nei Recommender Systems}

Nel contesto dei sistemi di raccomandazione, exploration ed exploitation presentano vantaggi e svantaggi differenti.

\subsubsection{Exploitation nei RecSys}

L'exploitation permette di:
\begin{itemize}
    \item raccomandare item che, sulla base delle conoscenze acquisite, risultano rilevanti per l'utente.
\end{itemize}

Tuttavia presenta diversi rischi:
\begin{itemize}
    \item la conoscenza delle preferenze dell'utente può restare parziale;
    \item il sistema può raccomandare item soltanto \emph{localmente ottimali}, cioè ottimali rispetto a ciò che è già noto;
    \item l'utente può essere progressivamente ``imprigionato in una bolla'', poiché il sistema continua a proporre contenuti simili a quelli già osservati.
\end{itemize}

\subsubsection{Exploration nei RecSys}

L'exploration permette invece di:
\begin{itemize}
    \item costruire una conoscenza più ampia delle preferenze dell'utente;
    \item aumentare la possibilità di ``stupire'' l'utente proponendo contenuti inattesi ma graditi;
    \item aprire l'utente a nuove esperienze.
\end{itemize}

L'esplorazione presenta però un costo:
\begin{itemize}
    \item può portare alla raccomandazione di item non rilevanti;
    \item raccomandazioni non pertinenti possono provocare una disaffezione dell'utente.
\end{itemize}

Il RecSys deve quindi posizionarsi opportunamente tra questi due estremi.

\subsection{Applicazioni del Multi-Armed Bandit}

Le tecniche MAB trovano applicazione in numerosi domini. Le slide citano, tra gli altri:
\begin{itemize}
    \item \emph{clinical trials};
    \item \emph{portfolio design};
    \item \emph{wireless selection};
    \item \emph{hyperparameter optimization};
    \item \emph{influence maximization};
    \item \emph{reinforcement learning}/Markov decision trees.
\end{itemize}

Come riferimento applicativo viene indicato:
\begin{quote}
D. Bouneffouf, I. Rish, \emph{A Survey on Practical Applications of Multi-Armed and Contextual Bandits}, 2019.
\end{quote}

La varietà dei contesti applicativi ha condotto alla formulazione di numerose varianti del problema MAB. Le slide ne considerano soltanto alcune.

\subsection{Arm stocastici}

Una delle assunzioni fondamentali più comuni nel Multi-Armed Bandit Problem è che l'esito di una scelta non sia deterministico.

Esempi:
\begin{itemize}
    \item in un \emph{clinical trial}, la somministrazione dello stesso farmaco non produce necessariamente lo stesso decorso in tutti i pazienti;
    \item l'andamento di un titolo finanziario non è deterministico;
    \item l'esito di una raccomandazione dipende dalla risposta dell'utente;
    \item l'esito della visualizzazione di un banner dipende dall'utente che lo osserva.
\end{itemize}

Nel linguaggio dei MAB:
\begin{itemize}
    \item ciascuna alternativa, azione o decisione prende il nome di \emph{arm};
    \item la selezione di un arm prende il nome di \emph{pull};
    \item ogni pull produce stocasticamente un \emph{reward}.
\end{itemize}

\subsubsection{Origine della terminologia ``Multi-Armed Bandit''}

Il nome deriva dall'analogia con le slot machine dei casinò.

Si immagini di avere a disposizione diverse slot machine. Ciascuna possiede un ``braccio'' (\emph{arm}) che può essere tirato (\emph{pull}). Ogni pull restituisce un certo reward.

La caratteristica fondamentale è che il comportamento della macchina è stocastico: a ciascuna slot machine è associata una propria distribuzione di probabilità sui reward.

Inizialmente tali distribuzioni sono sconosciute. Il giocatore deve quindi impararne le caratteristiche esclusivamente sulla base dei reward osservati dopo i pull effettuati.

\subsection{Protocollo formale del MAB stocastico}

Si consideri un insieme di $K$ arm:
\[
\mathcal{A}=\{1,\ldots,K\}.
\]

A ogni arm $i$ è associata una distribuzione di reward sconosciuta $D_i$. Nella notazione delle slide:
\[
D_i:\mathbb{R}\rightarrow[0,1].
\]

Il protocollo mostrato nelle slide assume:
\begin{itemize}
    \item $K$ arm;
    \item un orizzonte temporale di $T$ round, inizialmente considerato noto;
    \item una distribuzione di reward $D_a$ sconosciuta per ciascun arm $a$.
\end{itemize}

Per ciascun round $t\in\{1,\ldots,T\}$:
\begin{enumerate}
    \item l'algoritmo seleziona un arm $a_t$;
    \item viene estratto un reward
    \[
    r_t\in[0,1]
    \]
    dalla distribuzione associata all'arm selezionato:
    \[
    r_t\sim D_{a_t};
    \]
    \item il reward $r_t$ viene osservato dall'algoritmo.
\end{enumerate}

Il protocollo specifica inoltre che il reward è campionato indipendentemente dalla distribuzione dell'arm selezionato.

Un aspetto importante è che l'algoritmo osserva il reward dell'arm effettivamente selezionato e ``nulla altro'': non conosce cioè, nello stesso round, quale reward avrebbero prodotto gli arm non scelti.

\subsection{Esempio: banner e reward bernoulliani}

Si consideri nuovamente il problema di scegliere, tra $K$ differenti versioni di un banner pubblicitario, quella più adatta a una campagna.

I banner rappresentano gli arm:
\[
\mathcal{A}=\{1,\ldots,K\}.
\]

Se l'obiettivo è massimizzare il CTR, è naturale codificare il reward come:
\[
r=
\begin{cases}
1, & \text{se il banner viene cliccato},\\
0, & \text{se il banner non viene cliccato}.
\end{cases}
\]

A ogni banner $i$ è associata una probabilità $p_i$ di ricevere un click.

Pertanto la distribuzione $D_i$ è bernoulliana:
\[
D_i(r)=
\begin{cases}
p_i, & r=1,\\
1-p_i, & r=0.
\end{cases}
\]

In questo caso il reward atteso è:
\begin{align*}
\mu_i
&=\mathbb{E}[D_i]\\
&=1\cdot p_i+0\cdot(1-p_i)\\
&=p_i.
\end{align*}

Di conseguenza, nel caso dei banner bernoulliani, individuare l'arm con reward atteso massimo equivale a individuare il banner con probabilità di click massima.

\subsection{Obiettivo del MAB stocastico}

A ciascun arm $i$ è associato il valore atteso:
\[
\mu_i=\mathbb{E}[D_i].
\]

L'arm ottimale è indicato con:
\[
i^*=\arg\max_{i\in\mathcal{A}}\mu_i.
\]

Il suo reward atteso è:
\[
\mu^*=\max_{i\in\mathcal{A}}\mu_i.
\]

L'esecuzione di un algoritmo MAB produce una sequenza, o \emph{traccia},
\[
a_1,r_1,\ldots,a_t,r_t,\ldots,a_T,r_T,
\]
dove:
\begin{itemize}
    \item $a_t$ è l'arm selezionato al round $t$;
    \item $r_t$ è il reward prodotto da $D_{a_t}$.
\end{itemize}

L'obiettivo considerato nelle prime parti delle slide consiste nel massimizzare il \emph{cumulative reward}:
\[
R_T=\sum_{l=1}^{T}r_l.
\]

La strategia deve quindi acquisire informazioni sugli arm senza sacrificare eccessivamente il reward accumulato durante il processo.

\subsubsection{Strategie deterministiche}

Una strategia deterministica è una funzione che, data una traccia parziale
\[
a_1,r_1,\ldots,a_t,r_t,
\]
eventualmente vuota per $t=0$, restituisce un unico arm:
\[
a_{t+1}.
\]

Indicando la storia fino a $t$ con
\[
\mathcal{H}_t=(a_1,r_1,\ldots,a_t,r_t),
\]
una politica deterministica può quindi essere espressa astrattamente come
\[
\pi(\mathcal{H}_t)=a_{t+1}.
\]

A parità di storia, una strategia deterministica seleziona sempre lo stesso arm.

\subsubsection{Strategie stocastiche}

Una strategia stocastica, data la storia
\[
\mathcal{H}_t=(a_1,r_1,\ldots,a_t,r_t),
\]
restituisce invece una distribuzione di probabilità sugli arm.

Il prossimo arm viene campionato da tale distribuzione:
\[
a_{t+1}\sim \pi(\cdot\mid\mathcal{H}_t).
\]

Di conseguenza, anche a parità di storia, due esecuzioni differenti possono selezionare arm diversi.

\subsection{Notazione fondamentale}

Sia
\[
a_1,r_1,\ldots,a_t,r_t
\]
la traccia osservata fino al round $t$.

\subsubsection{Funzione caratteristica}

Le slide utilizzano la funzione caratteristica:
\[
\mathbb{I}(\mathrm{cond}),
\]
definita implicitamente da
\[
\mathbb{I}(\mathrm{cond})=
\begin{cases}
1, & \text{se la condizione è vera},\\
0, & \text{se la condizione è falsa}.
\end{cases}
\]

\subsubsection{Numero di pull di un arm}

Il numero di volte in cui l'arm $i$ è stato selezionato fino al round $t$ è:
\[
N_{i,t}
=
\sum_{h=1}^{t}\mathbb{I}(a_h=i).
\]

Ogni termine della somma vale $1$ soltanto nei round in cui è stato selezionato l'arm $i$.

\subsubsection{Media empirica}

La media empirica dei reward ottenuti dall'arm $i$ fino al round $t$ è:
\[
\widehat{\mu}_{i,t}
=
\frac{
\sum_{h=1}^{t}\mathbb{I}(a_h=i)r_h
}{
N_{i,t}
}.
\]

Il numeratore seleziona esclusivamente i reward osservati nei round in cui $a_h=i$, mentre il denominatore è il numero complessivo di osservazioni dell'arm.

La quantità $\widehat{\mu}_{i,t}$ rappresenta quindi la stima empirica di $\mu_i$.

\subsubsection{Deviazione standard empirica}

Le slide definiscono inoltre:
\[
\widehat{\sigma}_{i,t}
=
\sqrt{
\frac{
\sum_{h=1}^{t}
\mathbb{I}(a_h=i)
\left(r_h-\widehat{\mu}_{i,t}\right)^2
}{
N_{i,t}-1
}
}.
\]

Essa rappresenta la deviazione standard empirica dei reward osservati per l'arm $i$.

\subsection{Explore-Then-Commit (ETC)}

Una prima politica di allocazione è \emph{Explore-Then-Commit} (ETC).

L'idea consiste nel dividere l'esecuzione in due fasi nettamente separate:
\begin{enumerate}
    \item \textbf{exploration}: per una frazione $\gamma$ dell'orizzonte si esplorano uniformemente gli arm;
    \item \textbf{commit/exploitation}: terminata l'esplorazione, si sceglie l'arm con media empirica maggiore e lo si utilizza per il resto dell'orizzonte.
\end{enumerate}

Il parametro soddisfa:
\[
0\leq\gamma\leq1.
\]

La durata della fase esplorativa è:
\[
\gamma T.
\]

Nella notazione delle slide:
\[
a_{t+1}
=
\begin{cases}
(t+1)\bmod K,
& t\leq\gamma T,\\[1ex]
\displaystyle
\arg\max_{i\in\mathcal{A}}
\widehat{\mu}_{i,\gamma T},
& t>\gamma T.
\end{cases}
\]

L'espressione modulo $K$ rappresenta l'esplorazione ciclica, o \emph{round-robin}, dei diversi arm.

Questa politica viene associata nelle slide anche all'\emph{A/B testing}.

\subsubsection{Scelta del parametro $\gamma$}

L'obiettivo continua a essere la massimizzazione del cumulative reward $R_T$.

Se $\gamma$ è troppo piccolo:
\begin{itemize}
    \item la fase esplorativa raccoglie pochi campioni;
    \item le medie empiriche possono essere stime poco affidabili;
    \item al termine dell'esplorazione potrebbe essere scelto
    \[
    j=
    \arg\max_{i\in\mathcal{A}}
    \widehat{\mu}_{i,\gamma T}
    \]
    anche se $j$ è sub-ottimale;
    \item si potrebbe quindi avere
    \[
    \mu_j\ll\mu^*;
    \]
    \item poiché $j$ viene poi selezionato per tutti i round
    \[
    \gamma T<t\leq T,
    \]
    l'errore commesso alla fine dell'esplorazione viene propagato per tutto il resto dell'esecuzione.
\end{itemize}

Se invece $\gamma$ è troppo vicino a $1$:
\begin{itemize}
    \item quasi tutti i pull vengono dedicati all'esplorazione uniforme;
    \item gli arm vengono selezionati un numero pressoché uguale di volte;
    \item questo accade indipendentemente dai reward che stanno producendo;
    \item si verifica quindi un fenomeno di \emph{over-exploration}.
\end{itemize}

Un ulteriore limite riguarda la modalità di esplorazione: una distribuzione uniforme dei pull può essere inefficiente, perché tratta allo stesso modo arm promettenti e arm chiaramente poco promettenti.

Il problema di migliorare la fase di \emph{pure exploration} conduce alla branca denominata \emph{Best Arm Identification} (BAI), ripresa nelle slide finali.

\subsection{$\epsilon$-Greedy}

A differenza di ETC, $\epsilon$-Greedy non separa rigidamente exploration ed exploitation, ma le alterna durante l'intera esecuzione.

Dato:
\[
0\leq\epsilon\leq1,
\]
la politica opera come segue:
\begin{itemize}
    \item con probabilità $1-\epsilon$, effettua exploitation scegliendo l'arm con media empirica maggiore;
    \item con probabilità $\epsilon$, effettua exploration scegliendo casualmente un arm.
\end{itemize}

Si tratta quindi di una strategia stocastica.

\subsubsection{Inizializzazione}

Per i primi $K$ round viene effettuato un pull di ciascun arm:
\[
a_{t+1}=t+1,
\qquad t<K.
\]

L'obiettivo dell'inizializzazione è garantire che per ogni arm esista almeno un reward osservato, rendendo quindi definibili le medie empiriche.

\subsubsection{Politica di selezione}

Successivamente:
\begin{align*}
p &\sim \operatorname{rand}[0,1],\\
\text{se }p<\epsilon:\qquad
a_{t+1}&=\operatorname{rand}(K),\\
\text{altrimenti}:\qquad
a_{t+1}
&=
\arg\max_{i\in\mathcal{A}}\widehat{\mu}_{i,t}.
\end{align*}

La prima scelta è \emph{explorative}; la seconda è \emph{exploitive}.

\subsubsection{Casi estremi}

Per
\[
\epsilon=1,
\]
la politica seleziona sempre casualmente un arm. L'intera esecuzione è quindi orientata all'exploration e la conoscenza acquisita non viene utilizzata per fare exploitation.

Per
\[
\epsilon=0,
\]
la politica seleziona sempre un arm con media empirica massima. Il comportamento è completamente orientato all'exploitation.

Il significato di un valore intermedio di $\epsilon$ è precisamente quello di bilanciare i due comportamenti.

\subsubsection{Politiche Anytime}

ETC richiede di conoscere l'orizzonte $T$, perché la durata della fase esplorativa è definita tramite $\gamma T$.

$\epsilon$-Greedy non richiede invece di conoscere preventivamente $T$. Politiche di questo tipo vengono indicate nelle slide come \emph{AnyTime}.

Questa proprietà è particolarmente importante nelle applicazioni nelle quali non è naturale stabilire a priori il numero complessivo di interazioni, come nel caso dei sistemi di raccomandazione.

\subsection{Scelta di $\epsilon$}

Il valore di $\epsilon$ è un iperparametro.

Come accade per $\gamma$ in ETC, sorge quindi un problema di \emph{hyper-parameter optimization}: quale valore di $\epsilon$ produce il comportamento migliore?

Il valore ottimale dipende dalla particolare istanza del MAB:
\begin{itemize}
    \item numero di arm;
    \item valori attesi $\mu_i$;
    \item varianze delle distribuzioni;
    \item più in generale, caratteristiche delle distribuzioni associate agli arm.
\end{itemize}

Esiste tuttavia una difficoltà concettuale: se tali caratteristiche fossero completamente note, si conoscerebbe già quale arm possiede il valore atteso massimo e il problema di apprendimento sarebbe in larga parte risolto.

Le slide indicano:
\[
\epsilon=0.1
\]
come valore frequentemente utilizzato, sottolineando però che non è necessariamente ottimale per tutte le istanze.

Un'altra possibilità consiste nel considerare una classe di problemi, testare differenti valori di $\epsilon$ e scegliere quello che produce mediamente risultati migliori.

\subsection{Esperimenti con arm bernoulliani}

Le slide presentano un esperimento con:
\[
K=5
\]
arm bernoulliani e probabilità:
\[
p=\{0.1,0.1,0.1,0.1,0.9\}.
\]

L'ultimo arm è quindi nettamente migliore degli altri.

L'orizzonte utilizzato è:
\[
T=250.
\]

I valori testati sono:
\[
\epsilon\in\{0.1,0.2,0.3,0.4,0.5\}.
\]

Per ogni configurazione vengono eseguite:
\[
5000
\]
simulazioni.

I risultati visualizzati sono medie sui reward ottenuti nelle diverse simulazioni e rappresentano pertanto il comportamento medio.

Le slide mostrano due grafici:
\begin{itemize}
    \item la probabilità di selezionare il \emph{best arm} in funzione del tempo;
    \item il cumulative reward in funzione del tempo.
\end{itemize}

Nel secondo grafico della slide compare testualmente il titolo ``Cumulative Reward of the Exp3 Algorithm'' e la legenda ``Gamma'', pur trovandosi all'interno della discussione sperimentale dedicata ai differenti valori di $\epsilon$.

\subsection{Stabilizzazione di $\epsilon$-Greedy}

Un'importante osservazione sperimentale è che l'accuratezza di $\epsilon$-Greedy tende prima o poi a stabilizzarsi.

Inoltre:
\begin{itemize}
    \item valori più grandi di $\epsilon$ tendono a produrre una stabilizzazione più rapida;
    \item la stabilizzazione non è un fenomeno accidentale legato soltanto all'esperimento mostrato, ma deriva dalla struttura stessa dell'algoritmo.
\end{itemize}

\subsubsection{Probabilità di selezione di un arm}

Sia:
\[
h=
\arg\max_{i\in\mathcal{A}}
\widehat{\mu}_{i,t}
\]
l'arm che al round $t$ possiede la maggiore media empirica.

La probabilità di scegliere un generico arm $i$ è:
\[
\Pr(a_t=i)
=
\begin{cases}
(1-\epsilon)+\dfrac{\epsilon}{K},
& i=h,\\[2ex]
\dfrac{\epsilon}{K},
& i\neq h.
\end{cases}
\]

La prima espressione deriva da due possibilità:
\begin{enumerate}
    \item con probabilità $1-\epsilon$ viene effettuata exploitation e viene scelto sicuramente $h$;
    \item con probabilità $\epsilon$ viene effettuata exploration e, scegliendo uniformemente uno tra $K$ arm, $h$ viene comunque estratto con probabilità $1/K$.
\end{enumerate}

Pertanto:
\[
\Pr(a_t=h)
=
1-\epsilon+\frac{\epsilon}{K}.
\]

Per qualsiasi arm $i$ vale comunque:
\[
\Pr(a_t=i)\geq\frac{\epsilon}{K}.
\]

Se $\epsilon>0$, questa probabilità è strettamente positiva e non scompare nel tempo.

\subsubsection{Ogni arm viene campionato infinite volte}

Su un'esecuzione di durata infinita, le slide esprimono la conseguenza precedente mediante:
\[
\Pr\left(
\lim_{t\rightarrow\infty}N_{i,t}=\infty
\right)=1.
\]

In altre parole, con probabilità $1$ ogni arm viene selezionato infinite volte.

\subsubsection{Convergenza delle medie empiriche}

La legge dei grandi numeri fornisce:
\[
\Pr\left(
\lim_{N_{i,t}\rightarrow\infty}
\widehat{\mu}_{i,t}
=
\mu_i
\right)=1.
\]

Poiché per $\epsilon>0$ ogni arm viene campionato infinite volte, combinando i due risultati si ottiene:
\[
\Pr\left(
\lim_{t\rightarrow\infty}
\widehat{\mu}_{i,t}
=
\mu_i
\right)=1.
\]

Dunque le medie empiriche convergono quasi certamente ai rispettivi valori attesi.

Quando le stime sono sufficientemente vicine ai veri valori attesi, l'arm con media empirica massima coincide con l'arm ottimale:
\[
\arg\max_{i\in\mathcal{A}}
\widehat{\mu}_{i,t}
=
i^*.
\]

Da quel momento, salvo le scelte esplorative casuali, la componente greedy seleziona l'arm corretto.

\subsubsection{Accuratezza asintotica}

Una volta stabilizzatasi l'identificazione del best arm, la probabilità di selezionarlo è:
\[
(1-\epsilon)+\frac{\epsilon}{K}.
\]

Questa quantità rappresenta quindi il livello intorno al quale tende a stabilizzarsi l'accuratezza.

\subsection{Interpretazione dei test di $\epsilon$-Greedy}

Per:
\[
\epsilon\in\{0.2,0.3,0.4,0.5\},
\]
nel test mostrato l'accuratezza tende a stabilizzarsi intorno ai primi $100$ round.

Ciò indica che l'arm con media empirica maggiore è ormai, di fatto, l'arm ottimale.

Tuttavia, un grande $\epsilon$ implica una maggiore frequenza di pull casuali e quindi riduce la percentuale asintotica di selezioni ottimali.

\subsubsection{Caso $\epsilon=0.3$}

Con:
\[
\epsilon=0.3,
\]
la componente di exploitation viene utilizzata con probabilità:
\[
1-\epsilon=0.7.
\]

Una volta individuato il best arm, questa componente seleziona l'arm ottimale.

Nel restante $30\%$ dei casi viene scelto un arm casualmente. Poiché $K=5$, durante una scelta casuale il best arm viene selezionato con probabilità:
\[
\frac{1}{5}.
\]

Pertanto:
\[
0.7+\frac{0.3}{5}
=
0.7+0.06
=
0.76.
\]

L'accuratezza asintotica è dunque:
\[
0.76.
\]

\subsubsection{Caso $\epsilon=0.1$}

Per:
\[
\epsilon=0.1,
\]
la stabilizzazione è più lenta e diventa evidente soltanto verso la fine dei $250$ round.

Intorno al round $120$, tuttavia, la percentuale di scelte ottimali supera quella delle configurazioni con $\epsilon$ maggiore.

Il valore asintotico è:
\[
0.9+\frac{0.1}{5}
=
0.9+0.02
=
0.92.
\]

Quindi un $\epsilon$ più piccolo:
\begin{itemize}
    \item rallenta l'acquisizione dell'informazione necessaria a identificare l'arm migliore;
    \item una volta identificato l'arm migliore, permette però di selezionarlo con frequenza più elevata.
\end{itemize}

\subsection{Analisi del cumulative reward}

L'accuratezza descrive la probabilità puntuale di selezionare il best arm a un determinato istante.

L'obiettivo del MAB considerato nelle slide è però massimizzare:
\[
R_T=\sum_{t=1}^{T}r_t.
\]

Il cumulative reward dipende dall'intera traiettoria precedente e anche dall'orizzonte $T$.

Nel test mostrato:
\begin{itemize}
    \item per $T=250$, il valore di $\epsilon$ che produce il cumulative reward migliore è
    \[
    \epsilon=0.2;
    \]
    \item se l'orizzonte fosse soltanto
    \[
    T=150,
    \]
    il valore migliore sarebbe invece
    \[
    \epsilon=0.3.
    \]
\end{itemize}

A partire approssimativamente da:
\[
t=120,
\]
l'accuratezza di $\epsilon=0.1$ supera quella delle configurazioni con $\epsilon$ più alto.

Da quel momento il cumulative reward associato a $\epsilon=0.1$ cresce più rapidamente. Ne consegue che, per $T$ sufficientemente grande, $\epsilon=0.1$ finirà per produrre un cumulative reward superiore.

Le slide osservano inoltre che una strategia completamente exploitive:
\[
\epsilon=0
\]
avrebbe una curva di accuratezza ancora più lenta e potrebbe raggiungere il cumulative reward delle altre configurazioni soltanto dopo un numero considerevole di round.

\subsection{Annealing $\epsilon$-Greedy}

Gli esperimenti evidenziano due effetti:
\begin{itemize}
    \item un $\epsilon$ grande identifica rapidamente il best arm grazie alla maggiore esplorazione;
    \item un $\epsilon$ grande continua però a pagare il costo di numerosi pull casuali anche quando il best arm è ormai noto.
\end{itemize}

Da ciò nasce l'idea di utilizzare un valore di $\epsilon$ variabile nel tempo:
\begin{enumerate}
    \item inizialmente $\epsilon$ è elevato, favorendo exploration;
    \item progressivamente $\epsilon$ viene ridotto;
    \item la politica diventa quindi sempre più exploitive.
\end{enumerate}

Questa strategia è denominata \emph{Annealing $\epsilon$-Greedy}.

La politica è:
\begin{align*}
p &\sim \operatorname{rand}[0,1],\\
p<\epsilon_t
&\quad\Longrightarrow\quad
a_{t+1}=\operatorname{rand}(K),\\
p\geq\epsilon_t
&\quad\Longrightarrow\quad
a_{t+1}
=
\arg\max_{i\in\mathcal{A}}\widehat{\mu}_{i,t},\\
\epsilon_{t+1}&=f(t),
\end{align*}
dove:
\[
f(t)
\]
è una funzione decrescente di $t$.

\subsubsection{Funzioni di decadimento}

La qualità della politica dipende dalla scelta di $f$.

Una forma semplice è il decadimento esponenziale:
\[
f(t)=\beta e^{-\alpha t}.
\]

Le slide osservano tuttavia che questa scelta tende generalmente a ridurre $\epsilon$ troppo velocemente.

È utile invece mantenere la strategia sufficientemente esplorativa nei primi round.

Una funzione proposta nelle slide, utilizzata anche nel reinforcement learning, è:
\[
f(t)
=
\frac{1}
{1+\beta e^{\alpha(t-\gamma)}}.
\]

Poiché il denominatore cresce al crescere di $t$, $f(t)$ decresce e quindi la politica si sposta progressivamente dall'exploration verso l'exploitation.

\subsection{Limiti di $\epsilon$-Greedy}

Il principale limite di $\epsilon$-Greedy è che il livello di esplorazione è controllato da un fattore $\epsilon$ che, sia esso costante oppure funzione del tempo, non dipende direttamente dalle distribuzioni associate agli arm.

Si considerino due arm bernoulliani con probabilità:
\[
p_1,\qquad p_2
\]
e:
\[
\epsilon=0.1.
\]

\subsubsection{Arm difficili da distinguere}

Se:
\[
p_1\approx p_2,
\]
le due distribuzioni sono molto simili.

Sono quindi necessari molti campioni per distinguerle. Un'esplorazione del solo $10\%$ può risultare insufficiente:
\[
\epsilon=0.1
\]
è quindi potenzialmente troppo piccolo.

\subsubsection{Arm facilmente distinguibili}

Se $p_1$ e $p_2$ sono invece molto differenti, l'arm migliore può essere identificato con un numero relativamente piccolo di osservazioni.

Continuare a esplorare nel $10\%$ dei round può allora essere eccessivo:
\[
\epsilon=0.1
\]
può risultare troppo grande.

\subsubsection{Esplorazione ``alla cieca''}

Quando $\epsilon$-Greedy decide di esplorare, seleziona un arm casualmente senza distinguere tra:
\begin{itemize}
    \item arm apparentemente promettenti;
    \item arm apparentemente molto scarsi.
\end{itemize}

Le slide suggeriscono che una strategia più raffinata dovrebbe assegnare maggiore attenzione agli arm che appaiono più promettenti.

Da questa osservazione derivano strategie nelle quali si costruisce esplicitamente una misura della ``promettentezza'' di ciascun arm.

\subsection{Softmax}

La stima più immediata della qualità dell'arm $i$ al round $t$ è:
\[
\widehat{\mu}_{i,t}.
\]

Una prima idea consiste nell'assegnare una probabilità di selezione direttamente proporzionale alla media empirica:
\[
\mathcal{P}(i,t)
=
\frac{
\widehat{\mu}_{i,t}
}{
\sum_{j=1}^{K}\widehat{\mu}_{j,t}
}.
\]

\subsubsection{Perché la proporzionalità diretta è troppo esplorativa}

Per arm bernoulliani non degeneri, al crescere di $t$ si ha, quasi sicuramente:
\[
\mathcal{P}(i,t)>0
\]
per ogni arm.

Di conseguenza tutti gli arm continuano a essere selezionati e, nel limite, vengono osservati un numero infinito di volte.

Per la legge dei grandi numeri:
\[
\widehat{\mu}_{i,t}\rightarrow\mu_i.
\]

La probabilità di selezione tende quindi a:
\[
\mathcal{P}(i,t)
\longrightarrow
\frac{
\mu_i
}{
\sum_{j=1}^{K}\mu_j
}.
\]

Se:
\begin{itemize}
    \item il numero di arm è elevato, oppure
    \item i reward attesi sono relativamente simili,
\end{itemize}
allora:
\[
\sum_{j=1}^{K}\mu_j
\]
può essere molto maggiore di ciascun singolo $\mu_i$.

Anche il best arm può quindi mantenere una probabilità relativamente bassa di essere selezionato. Questo penalizza l'accuratezza.

\subsubsection{Distribuzione di Boltzmann o Softmax}

Per accentuare maggiormente le differenze tra le medie si utilizza una distribuzione di Boltzmann:
\[
\mathcal{P}(i,t)
=
\frac{
e^{\alpha\widehat{\mu}_{i,t}}
}{
\sum_{j=1}^{K}
e^{\alpha\widehat{\mu}_{j,t}}
}.
\]

Questa trasformazione è denominata \emph{softmax}.

La probabilità è correttamente normalizzata:
\[
\sum_{i=1}^{K}\mathcal{P}(i,t)=1.
\]

Per ragioni storiche, l'inverso:
\[
\tau=\frac{1}{\alpha}
\]
viene chiamato \emph{temperatura} della distribuzione.

\subsubsection{Rapporto tra le probabilità di due arm}

Consideriamo due arm $i$ e $j$.

Si ha:
\begin{align*}
\frac{\mathcal{P}(i,t)}
{\mathcal{P}(j,t)}
&=
\frac{
e^{\alpha\widehat{\mu}_{i,t}}
}{
e^{\alpha\widehat{\mu}_{j,t}}
}\\
&=
e^{\alpha
(\widehat{\mu}_{i,t}-\widehat{\mu}_{j,t})}.
\end{align*}

Definendo:
\[
\Delta_{ij}
=
\widehat{\mu}_{i,t}
-
\widehat{\mu}_{j,t},
\]
otteniamo:
\[
\frac{\mathcal{P}(i,t)}
{\mathcal{P}(j,t)}
=
e^{\alpha\Delta_{ij}}.
\]

Questa relazione rende immediato il significato del parametro $\alpha$.

\paragraph{Caso $\alpha\rightarrow0$.}
Per:
\[
\alpha\rightarrow0,
\]
vale:
\[
e^{\alpha\Delta_{ij}}\rightarrow1.
\]

Le probabilità dei differenti arm diventano quindi uguali e ciascun arm tende ad avere probabilità:
\[
\frac{1}{K}.
\]

La strategia è fortemente explorative.

\paragraph{Caso $\alpha\rightarrow\infty$.}
Per valori molto grandi di $\alpha$, anche piccole differenze $\Delta_{ij}$ vengono amplificate esponenzialmente.

L'arm con media empirica maggiore riceve quindi una probabilità molto superiore alle altre.

La strategia diventa fortemente exploitive.

\paragraph{Dipendenza dalla difficoltà dell'istanza.}
A $\alpha$ fissato, se:
\[
|\Delta_{ij}|
\]
è piccolo, le probabilità degli arm restano simili e la strategia esplora maggiormente.

Se invece le differenze tra le medie empiriche sono grandi, Softmax concentra naturalmente più probabilità sugli arm migliori.

Questa proprietà differenzia Softmax da $\epsilon$-Greedy: la tendenza a esplorare non dipende soltanto da un parametro prefissato, ma anche da quanto le alternative risultano effettivamente distinguibili.

La slide grafica confronta:
\[
\alpha=1,\qquad
\alpha=2,\qquad
\alpha=0.5,
\]
mostrando la funzione:
\[
e^{\alpha\Delta_{ij}}.
\]

La crescita è:
\begin{itemize}
    \item più accentuata per $\alpha=2$;
    \item intermedia per $\alpha=1$;
    \item più graduale per $\alpha=0.5$.
\end{itemize}

\subsection{Annealing Softmax}

Anche Softmax può essere resa maggiormente esplorativa all'inizio e progressivamente più exploitive.

Poiché valori piccoli di $\alpha$ favoriscono exploration e valori grandi favoriscono exploitation, occorre scegliere:
\[
\alpha(t)
\]
crescente con il tempo.

Le slide propongono come esempio:
\[
\alpha(t)=c+\log t,
\qquad
0<c\ll1.
\]

All'aumentare di $t$:
\[
\alpha(t)\uparrow,
\]
e quindi la distribuzione Softmax tende progressivamente a concentrarsi maggiormente sugli arm con media empirica più elevata.

Equivalentemente, poiché:
\[
\tau(t)=\frac{1}{\alpha(t)},
\]
la temperatura diminuisce nel tempo.

\subsection{Upper Confidence Bound (UCB)}

Le slide introducono UCB partendo dai limiti comuni di $\epsilon$-Greedy e Softmax.

Entrambe:
\begin{itemize}
    \item sono strategie stocastiche;
    \item tendono a preferire l'arm con media empirica maggiore;
    \item occasionalmente esplorano arm differenti.
\end{itemize}

La differenza è che:
\begin{itemize}
    \item $\epsilon$-Greedy, in exploration, seleziona un arm casualmente;
    \item Softmax attribuisce maggiore probabilità agli arm con media empirica più alta.
\end{itemize}

In entrambi i casi la scelta dipende in maniera sostanziale dalla stima:
\[
\widehat{\mu}_{i,t}.
\]

Il problema è che le due strategie considerano principalmente \emph{ciò che si sa}, senza rappresentare esplicitamente \emph{quanto bene lo si sa}.

In particolare, non tengono direttamente conto dell'incertezza associata alla stima.

\subsubsection{Problema delle false partenze}

Si consideri un arm realmente ottimale che, per effetto della casualità, produca reward bassi nei primi pull.

La sua media empirica:
\[
\widehat{\mu}_{i,t}
\]
può inizialmente essere molto bassa.

Una strategia che utilizza soltanto la media può quindi smettere di selezionarlo frequentemente.

Serviranno potenzialmente molti round prima che nuove osservazioni riescano a ``invertire'' le medie empiriche e correggere la falsa impressione iniziale.

UCB affronta precisamente questo problema introducendo una misura dell'incertezza.

\subsubsection{Principio generale}

UCB associa a ciascun arm un valore:
\[
UCB(\mathcal{H},i)
\]
che combina:
\begin{itemize}
    \item una stima del reward;
    \item una componente legata all'incertezza.
\end{itemize}

La politica è deterministica:
\[
a_{t+1}
=
\arg\max_{i\in\mathcal{A}}
UCB(\mathcal{H},i),
\]
dove:
\[
\mathcal{H}
=
a_1,r_1,\ldots,a_t,r_t
\]
rappresenta la storia.

\subsubsection{Tre varianti riportate nelle slide}

Una prima variante è:
\[
UCB(\mathcal{H},i)
=
\widehat{\mu}_{i,t}
+
\sqrt{
\frac{
2\log(1/\delta)
}{
N_{i,t}
}
},
\qquad
0<\delta<1.
\]

Una seconda variante è:
\[
UCB(\mathcal{H},i)
=
\widehat{\mu}_{i,t}
+
\sqrt{
\frac{
2\log t
}{
N_{i,t}
}
}.
\]

Una terza variante introduce un parametro moltiplicativo $\alpha$:
\[
UCB(\mathcal{H},i)
=
\widehat{\mu}_{i,t}
+
\alpha
\sqrt{
\frac{
2\log t
}{
N_{i,t}
}
},
\qquad
\alpha>0.
\]

Se:
\[
N_{i,t}=0,
\]
le slide inizializzano:
\[
UCB(\mathcal{H},i)=+\infty.
\]

Questo garantisce che ciascun arm venga selezionato almeno una volta.

\subsection{Componente exploitive ed explorative di UCB}

La struttura generale è:
\[
UCB(\mathcal{H},i)
=
\underbrace{\widehat{\mu}_{i,t}}_{\text{exploitive}}
+
\underbrace{\mathrm{explorative\_bonus}}_{\text{explorative}}.
\]

La media empirica premia gli arm che hanno prodotto reward elevati.

Il bonus esplorativo premia invece gli arm sui quali vi è maggiore incertezza.

\subsubsection{Dipendenza da $N_{i,t}$}

Il bonus è decrescente rispetto a:
\[
N_{i,t}.
\]

L'intuizione è:
\begin{itemize}
    \item se $N_{i,t}$ è piccolo, l'arm è stato osservato poche volte e la stima è incerta;
    \item se $N_{i,t}$ è grande, l'arm è stato osservato molte volte e vi è meno necessità di esplorarlo.
\end{itemize}

Per le versioni contenenti $\log t$, le slide mostrano la relazione:
\[
\frac{\log t}{N_{i,t-1}+1}
<
\frac{\log(t-1)}{N_{i,t-1}}
<
\frac{\log t}{N_{i,t-1}}.
\]

La relazione evidenzia due situazioni.

\paragraph{L'arm $i$ viene selezionato.}
Il contatore passa da:
\[
N_{i,t-1}
\]
a:
\[
N_{i,t-1}+1,
\]
e il bonus esplorativo associato all'arm tende a diminuire.

\paragraph{Viene selezionato un altro arm.}
Il contatore $N_{i,t}$ rimane invariato, mentre il tempo globale aumenta da $t-1$ a $t$.

Il termine $\log t$ cresce, per cui il bonus associato all'arm non selezionato aumenta.

Questo meccanismo evita che un arm venga ignorato per sempre: se non viene scelto per molto tempo, il suo bonus di esplorazione tende a diventare progressivamente più importante.

Nella prima variante il bonus dipende essenzialmente dal numero di volte in cui $i$ è stato selezionato.

Nella seconda e nella terza variante compare anche il tempo globale $t$, e quindi il valore dipende indirettamente anche da quante volte sono stati selezionati gli altri arm.

\subsection{Ruolo del parametro $\alpha$ in UCB}

Consideriamo:
\[
UCB(\mathcal{H},i)
=
\widehat{\mu}_{i,t}
+
\sqrt{
\frac{2\log t}{N_{i,t}}
}.
\]

La componente esplorativa dipende esclusivamente da:
\begin{itemize}
    \item $N_{i,t}$;
    \item il numero complessivo di pull $t$.
\end{itemize}

La componente exploitive:
\[
\widehat{\mu}_{i,t}
\]
dipende invece direttamente dai reward e quindi dalla scala delle distribuzioni degli arm.

Le slide osservano che:
\begin{itemize}
    \item se i valori $\mu_i$ sono grandi, il bonus esplorativo può diventare relativamente trascurabile;
    \item se i valori $\mu_i$ sono prossimi a zero, il bonus può avere un peso relativamente eccessivo.
\end{itemize}

Da ciò nasce l'introduzione di:
\[
\alpha>0,
\]
che modula il peso relativo della componente esplorativa:
\[
UCB(\mathcal{H},i)
=
\widehat{\mu}_{i,t}
+
\alpha
\sqrt{
\frac{2\log t}{N_{i,t}}
}.
\]

\subsection{Optimism in the Face of Uncertainty}

Il principio alla base di UCB viene indicato come:
\begin{quote}
\emph{optimism in the face of uncertainty}.
\end{quote}

L'idea è costruire, intorno alla media empirica:
\[
\widehat{\mu}_{i,t},
\]
un intervallo:
\[
I_c
=
\left[
\widehat{\mu}_{i,t}-\gamma_{i,t},
\widehat{\mu}_{i,t}+\gamma_{i,t}
\right],
\]
nel quale il vero valore:
\[
\mu_i
\]
ricade con elevata probabilità.

Il parametro:
\[
\gamma_{i,t}
\]
rappresenta quindi l'ampiezza dell'incertezza.

Nella prima variante di UCB:
\[
\gamma_{i,t}
=
\sqrt{
\frac{
2\log(1/\delta)
}{
N_{i,t}
}
}.
\]

\subsubsection{Interpretazione della confidence bound}

Le slide riportano, per distribuzioni normali indicate come:
\[
D_i=\mathcal{N}(0,1),
\]
la relazione:
\[
\Pr\left(
\left|
\widehat{\mu}_{i,t}-\mu_i
\right|
\leq
\sqrt{
\frac{
2\log(1/\delta)
}{
N_{i,t}
}
}
\right)
\geq
1-2\delta.
\]

Questa relazione dà un'interpretazione statistica a $\gamma_{i,t}$: con probabilità almeno $1-2\delta$, l'errore tra media empirica e valore atteso è limitato dalla quantità indicata.

\subsubsection{Perché UCB è ``ottimista''}

Dato l'intervallo:
\[
I_c=
[
\widehat{\mu}_{i,t}-\gamma_{i,t},
\widehat{\mu}_{i,t}+\gamma_{i,t}
],
\]
UCB utilizza il valore più favorevole compatibile con tale intervallo:
\[
\widehat{\mu}_{i,t}+\gamma_{i,t}.
\]

Un arm molto incerto viene quindi valutato in maniera ``ottimistica'': finché i dati disponibili non permettono di escludere che possa essere molto buono, riceve un bonus che incentiva ulteriori osservazioni.

\subsubsection{Raffinamento della confidence}

Le prestazioni di UCB dipendono fortemente dalla qualità della stima di $\gamma_{i,t}$.

Le slide indicano due informazioni che possono essere sfruttate:
\begin{enumerate}
    \item la varianza delle distribuzioni $D_i$, stimabile empiricamente mediante
    \[
    \widehat{\sigma}_{i,t};
    \]
    \item la conoscenza della famiglia di distribuzioni, per esempio bernoulliane o normali.
\end{enumerate}

Da queste idee derivano varianti quali:
\begin{itemize}
    \item UCB-V;
    \item UCB-normal;
    \item KL-UCB.
\end{itemize}

\subsubsection{Confronto grafico delle strategie}

Le slide riportano inoltre un confronto grafico tra:
\begin{itemize}
    \item Annealing $\epsilon$-Greedy;
    \item UCB1;
    \item Annealing Softmax.
\end{itemize}

Il confronto viene effettuato sia in termini di:
\begin{itemize}
    \item probabilità di selezionare il best arm;
    \item cumulative reward.
\end{itemize}

Il grafico rende visibili differenti dinamiche transitorie nella fase iniziale e una progressiva crescita delle prestazioni all'aumentare del numero di round.

\subsection{Multi-Armed Bandit e sistemi di raccomandazione}

Nel contesto di un RecSys:
\begin{itemize}
    \item ogni arm corrisponde a un item;
    \item ogni step $t$ corrisponde a una richiesta/interazione di un utente;
    \item l'arm selezionato è l'item raccomandato all'utente.
\end{itemize}

\subsubsection{Insieme di arm variabile nel tempo}

Fino a questo punto è stato considerato prevalentemente un insieme fisso di $K$ arm.

Le slide osservano però che le tecniche possono essere adattate anche a un insieme di arm che varia nel tempo.

Indichiamo con:
\[
\mathcal{A}_t
\]
l'insieme degli arm disponibili allo step $t$.

Per UCB, un nuovo arm può essere inizializzato con:
\[
UCB(\mathcal{H},i)=+\infty.
\]

In questo modo l'arm appena introdotto viene esplorato.

La selezione diventa:
\[
a_{t+1}
=
\arg\max_{i\in\mathcal{A}_t}
UCB(\mathcal{H},i).
\]

Questa caratteristica rende i MAB particolarmente adatti ai contesti in cui utenti e item cambiano rapidamente.

Un esempio indicato nelle slide è quello dei portali di news, nei quali gli item sono articoli che vengono continuamente pubblicati e sostituiti.

\subsection{Raccomandazioni personalizzate e limite del MAB classico}

Per raccomandazioni non personalizzate, le strategie MAB possono essere utilizzate in modo abbastanza diretto.

Per esempio, una strategia \emph{Top ranked} può essere arricchita introducendo exploration, aumentando anche la diversificazione delle proposte.

La situazione cambia con le raccomandazioni personalizzate.

Supponiamo di utilizzare il click come reward:
\[
r=
\begin{cases}
1, & \text{item cliccato},\\
0, & \text{item non cliccato}.
\end{cases}
\]

Nel modello MAB classico a ogni item $i$ verrebbe associata una distribuzione bernoulliana:
\[
D_i.
\]

Questa distribuzione rappresenterebbe una probabilità generale di click dell'item.

Ma un RecSys personalizzato non è interessato soltanto alla probabilità che un item venga cliccato dalla popolazione generale.

Al round $t$ interessa invece sapere quanto lo specifico utente:
\[
u_t
\]
sia propenso a cliccare lo specifico item.

La risposta dipende quindi contemporaneamente:
\begin{itemize}
    \item dall'item;
    \item dall'utente.
\end{itemize}

Una singola distribuzione $D_i$ indipendente dall'utente non è più sufficiente.

\subsection{Contextual Bandits}

Per modellare situazioni come la raccomandazione personalizzata si introducono i \emph{Contextual Multi-Armed Bandit Problems}.

L'idea fondamentale è che la distribuzione di reward di un arm possa cambiare in funzione del contesto.

Al round $t$ compare un contesto:
\[
c_t.
\]

La distribuzione associata all'arm $i$ diventa quindi:
\[
D_{i,c_t}.
\]

Il contesto rappresenta informazione aggiuntiva che la strategia può utilizzare per effettuare la propria scelta.

Nei sistemi di raccomandazione, $c_t$ è tipicamente descritto mediante feature relative sia agli utenti sia agli item.

\subsubsection{Feature dell'utente}

Le slide riportano esempi quali:
\begin{itemize}
    \item età;
    \item provenienza;
    \item genere;
    \item interazioni passate.
\end{itemize}

\subsubsection{Feature degli item}

Nel caso di articoli di news:
\begin{itemize}
    \item tipologia dell'articolo, ad esempio cronaca nera, politica, sport;
    \item data di pubblicazione;
    \item autore.
\end{itemize}

Il comportamento dell'algoritmo può quindi dipendere dall'interazione tra le caratteristiche dell'utente corrente e quelle degli item disponibili.

\subsection{Linear Bandits}

I \emph{Linear Bandits} costituiscono un particolare modello contestuale.

Gli arm disponibili possono cambiare nel tempo:
\[
\mathcal{A}_t.
\]

Per ciascun arm:
\[
a\in\mathcal{A}_t
\]
e ciascun round $t$, il contesto genera un vettore di feature:
\[
\mathbf{x}_{t,a}\in\mathbb{R}^{d}.
\]

Il vettore è $d$-dimensionale:
\[
\mathbf{x}_{t,a}
=
(x_{t,a}^{(1)},\ldots,x_{t,a}^{(d)}),
\]
e rappresenta numericamente $d$ caratteristiche rilevanti.

I valori delle feature dipendono:
\begin{itemize}
    \item dall'arm/item $a$;
    \item dall'utente corrente $u_t$.
\end{itemize}

\subsubsection{Modello lineare del reward}

Le slide assumono un reward composto da:
\begin{itemize}
    \item una componente lineare;
    \item un rumore stocastico.
\end{itemize}

La relazione riportata è:
\[
D_{a,t}
=
\boldsymbol{\theta}^{*}
\circ
\mathbf{x}_{t,a}
+
\eta_t.
\]

Qui:
\begin{itemize}
    \item $\boldsymbol{\theta}^{*}\in\mathbb{R}^{d}$ è un vettore incognito;
    \item $\mathbf{x}_{t,a}\in\mathbb{R}^{d}$ è il vettore di feature;
    \item $\eta_t$ rappresenta il rumore stocastico;
    \item il simbolo $\circ$ è la notazione utilizzata nelle slide per il prodotto coinvolto nel modello lineare.
\end{itemize}

Tipicamente:
\[
\eta_t\sim\mathcal{N}(0,\sigma).
\]

Nel contesto dei RecSys, il rumore rappresenta possibili incertezze nel comportamento dell'utente.

Poiché il rumore ha media nulla, il reward atteso è:
\[
\mu_{a,t}
=
\mathbb{E}[D_{a,t}]
=
\boldsymbol{\theta}^{*}
\circ
\mathbf{x}_{t,a}.
\]

Il vettore sconosciuto $\boldsymbol{\theta}^{*}$ determina quindi il modo in cui le feature influenzano il reward atteso.

\subsection{UCB per Linear Bandits}

Il principio generale di UCB rimane invariato.

La selezione deve massimizzare una stima ottimistica del reward:
\[
a_t
=
\arg\max_{a\in\mathcal{A}_t}
ucb(a).
\]

Poiché:
\[
\boldsymbol{\theta}^{*}
\]
è sconosciuto, non è possibile calcolare direttamente:
\[
\boldsymbol{\theta}^{*}
\circ
\mathbf{x}_{t,a}.
\]

Occorre quindi rappresentare l'incertezza su $\boldsymbol{\theta}^{*}$.

\subsubsection{Confidence set $\Theta_t$}

Si introduce un insieme:
\[
\Theta_t\subset\mathbb{R}^{d}
\]
che deve soddisfare due esigenze:
\begin{enumerate}
    \item con buona probabilità:
    \[
    \boldsymbol{\theta}^{*}\in\Theta_t;
    \]
    \item $\Theta_t$ deve essere il più piccolo possibile.
\end{enumerate}

Il primo requisito garantisce un buon livello di confidenza.

Il secondo evita che l'incertezza sia talmente ampia da rendere poco informativa la stima.

Applicando il principio di \emph{optimism in the face of uncertainty}:
\[
ucb(a)
=
\max_{\boldsymbol{\theta}\in\Theta_t}
\boldsymbol{\theta}
\circ
\mathbf{x}_{t,a}.
\]

Per ciascun arm si considera quindi il reward migliore che esso potrebbe avere sotto un parametro compatibile con il confidence set corrente.

\subsection{Stima di $\boldsymbol{\theta}^{*}$}

Prima di determinare $\Theta_t$, occorre trovare una stima centrale:
\[
\widehat{\boldsymbol{\theta}}_t.
\]

Le slide costruiscono tale stima minimizzando una somma di errori quadratici più un termine di regolarizzazione:
\[
\widehat{\boldsymbol{\theta}}_t
=
\arg\min_{\boldsymbol{\theta}\in\mathbb{R}^{d}}
\left[
\sum_{s=1}^{t}
\left(
r_s-
\boldsymbol{\theta}
\circ
\mathbf{x}_{s,a_s}
\right)^2
+
\lambda
\|\boldsymbol{\theta}\|_2^2
\right].
\]

I termini hanno il seguente significato:
\begin{itemize}
    \item $r_s$ è il reward osservato al round $s$;
    \item $\mathbf{x}_{s,a_s}$ è il vettore di feature dell'arm effettivamente selezionato al round $s$;
    \item $\boldsymbol{\theta}\circ\mathbf{x}_{s,a_s}$ è la previsione lineare;
    \item
    \[
    r_s-\boldsymbol{\theta}\circ\mathbf{x}_{s,a_s}
    \]
    è l'errore di previsione;
    \item la somma degli errori quadratici misura quanto bene $\boldsymbol{\theta}$ spiega la storia osservata;
    \item
    \[
    \lambda\|\boldsymbol{\theta}\|_2^2
    \]
    è la componente di regolarizzazione;
    \item $\lambda$ è un iperparametro.
\end{itemize}

Le slide affermano che, differenziando l'obiettivo, si ottiene:
\[
\widehat{\boldsymbol{\theta}}_t
=
\left(
\lambda I
+
\sum_{s=1}^{t}
\mathbf{x}_{s,a_s}
\circ
\mathbf{x}_{s,a_s}
\right)^{-1}
\left(
\sum_{s=1}^{t}
r_s\mathbf{x}_{s,a_s}
\right),
\]
dove:
\[
I
\]
è la matrice identità di dimensione $d$.

Si osservi che, nella notazione delle slide, il simbolo $\circ$ viene utilizzato anche nel termine che deve essere compatibile dimensionalmente con la matrice $d\times d$ $\lambda I$; si mantiene qui la notazione originale del materiale.

\subsection{LinUCB e geometria dell'incertezza}

Una volta determinata:
\[
\widehat{\boldsymbol{\theta}}_t,
\]
si costruisce un intorno all'interno del quale si ritiene che:
\[
\boldsymbol{\theta}^{*}
\]
cada con elevata probabilità.

Le slide descrivono $\Theta_t$ intuitivamente come un ellissoide con centro:
\[
\widehat{\boldsymbol{\theta}}_t.
\]

Si definisce:
\[
V_t
=
\lambda I
+
\sum_{s=1}^{t}
\mathbf{x}_{s,a_s}
\circ
\mathbf{x}_{s,a_s}.
\]

Secondo la descrizione riportata nelle slide:
\begin{itemize}
    \item gli assi dell'ellissoide sono gli autovettori di $V_t$;
    \item le loro lunghezze sono gli inversi dei corrispondenti autovalori.
\end{itemize}

LinUCB considera l'incertezza lungo la direzione:
\[
\mathbf{x}_{t,a}.
\]

La politica presentata è:
\[
a_t
=
\arg\max_{a\in\mathcal{A}_t}
\left[
\widehat{\boldsymbol{\theta}}_t
\circ
\mathbf{x}_{t,a}
+
\alpha
\sqrt{
\mathbf{x}_{t,a}^{T}
V_t^{-1}
\mathbf{x}_{t,a}
}
\right].
\]

La funzione è nuovamente composta da due termini.

\paragraph{Componente exploitive.}
\[
\widehat{\boldsymbol{\theta}}_t
\circ
\mathbf{x}_{t,a}.
\]

È il reward previsto utilizzando la migliore stima corrente del parametro.

\paragraph{Componente explorative.}
\[
\alpha
\sqrt{
\mathbf{x}_{t,a}^{T}
V_t^{-1}
\mathbf{x}_{t,a}
}.
\]

Rappresenta l'incertezza lungo il vettore di feature dell'arm.

Il parametro:
\[
\alpha
\]
è un iperparametro che modula l'entità dell'\emph{explorative bonus}.

\subsection{Aggiornamento incrementale di LinUCB}

Le slide introducono gli aggiornamenti:
\[
V_{t+1}
=
V_t
+
\mathbf{x}_{t,a_t}
\circ
\mathbf{x}_{t,a_t},
\]
e:
\[
\mathbf{b}_{t+1}
=
\mathbf{b}_t
+
r_t\mathbf{x}_{t,a_t},
\]
dove:
\[
\mathbf{b}_t
=
\sum_{s=1}^{t}
r_s\mathbf{x}_{s,a_s}.
\]

La formula precedente per la stima può quindi essere riscritta utilizzando:
\[
V_t
=
\lambda I
+
\sum_{s=1}^{t}
\mathbf{x}_{s,a_s}
\circ
\mathbf{x}_{s,a_s}
\]
e:
\[
\mathbf{b}_t
=
\sum_{s=1}^{t}
r_s\mathbf{x}_{s,a_s}.
\]

In questa notazione:
\[
\widehat{\boldsymbol{\theta}}_t
=
V_t^{-1}\mathbf{b}_t.
\]

Dopo ogni nuova osservazione vengono quindi aggiornate le quantità che riassumono l'informazione raccolta fino a quel momento.

\subsection{Best Arm Identification}

Le slide finali introducono il \emph{Best Arm Identification} (BAI), in particolare nella variante \emph{fixed budget}.

Il modello di interazione resta quello di un MAB:
\begin{itemize}
    \item esiste un numero finito di arm stocastici;
    \item un \emph{forecaster} interagisce sequenzialmente con gli arm;
    \item a ogni round seleziona un arm;
    \item osserva un sample/reward estratto dalla distribuzione dell'arm selezionato.
\end{itemize}

Ciò che cambia è l'obiettivo.

\subsubsection{Differenza tra MAB classico e BAI}

Nel MAB discusso in precedenza si cerca di massimizzare:
\[
R_T=
\sum_{t=1}^{T}r_t.
\]

Quindi i reward ottenuti \emph{durante} il processo contano direttamente.

Nel Best Arm Identification, invece, non è necessario massimizzare il cumulative reward durante la fase di interazione.

L'obiettivo è, dopo un numero prefissato di round, identificare correttamente l'arm con massimo expected reward:
\[
i^*
=
\arg\max_i\mu_i.
\]

Di conseguenza, può essere accettabile sacrificare molti reward durante l'esplorazione purché ciò permetta di identificare con maggiore affidabilità l'arm migliore.

\subsubsection{Contesti applicativi}

BAI è indicato per situazioni in cui esiste una fase di testing offline.

Durante tale fase:
\begin{itemize}
    \item non è prioritario accumulare reward;
    \item è prioritario esplorare sufficientemente il comportamento dei differenti arm;
    \item al termine viene identificato il migliore;
    \item il best arm viene quindi impiegato stabilmente quando il sistema entra online.
\end{itemize}

Le slide citano:
\begin{itemize}
    \item Explore-Then-Commit;
    \item tuning di parametri discreti.
\end{itemize}

\subsection{Algoritmo generico per Fixed-Budget Best Arm Identification}

In termini astratti, la struttura di un algoritmo per Fixed-Budget BAI è simile a quella di un algoritmo MAB.

Si dispone di un budget:
\[
T,
\]
cioè del numero complessivo di round disponibili.

Finché il budget non è esaurito:
\begin{enumerate}
    \item si seleziona un arm $a_{t+1}$;
    \item si esegue un pull;
    \item si osserva un reward, indicato nella slide con $x_{t+1}$;
    \item si aggiunge alla storia la coppia:
    \[
    (a_{t+1},x_{t+1}).
    \]
\end{enumerate}

Si noti che in questa parte delle slide il simbolo $x_{t+1}$ viene utilizzato per il reward osservato e non per il vettore di contesto $\mathbf{x}_{t,a}$ utilizzato precedentemente nei Linear Bandits.

Al termine del budget viene restituito l'arm con media empirica più alta:
\[
\widehat{i}
=
\arg\max_{i\in\mathcal{A}}
\widehat{\mu}_{i,T}.
\]

Lo pseudocodice mostrato nelle slide può essere rappresentato concettualmente come:
\begin{enumerate}
    \item input: insieme di arm e budget $T$;
    \item per ogni $t<T$:
    \begin{enumerate}
        \item $a_{t+1}\leftarrow \mathrm{select}(\pi_t,T)$;
        \item $x_{t+1}\leftarrow\mathrm{sample\ arm}\ a_{t+1}$;
        \item aggiornare la storia con $(a_{t+1},x_{t+1})$;
    \end{enumerate}
    \item restituire:
    \[
    \arg\max_{i\in\mathcal{A}}\widehat{\mu}_{i,T}.
    \]
\end{enumerate}

Il problema fondamentale consiste quindi nel definire una buona strategia di selezione durante i $T$ round.

\subsection{Uniform Exploration nel Fixed-Budget BAI}

La strategia più semplice, e indicata nelle slide come meno efficace, è l'\emph{uniform exploration}, detta anche \emph{round-robin}.

La selezione ciclica viene espressa come:
\[
a_{t+1}
\leftarrow
(t+1)\bmod K.
\]

Tutti gli arm ricevono quindi approssimativamente lo stesso numero di sample.

\subsubsection{Limite della strategia uniforme}

La strategia è \emph{memoryless}: l'allocazione dei pull non cambia in funzione dell'andamento osservato round dopo round.

Le slide osservano che essa non dipende da $\pi_t$; il simbolo $\pi_t$ non viene ulteriormente formalizzato nel materiale, ma il punto espresso è che l'allocazione uniforme non si adatta all'informazione raccolta nel corso dell'esecuzione.

In particolare:
\begin{itemize}
    \item un arm chiaramente scarso continua a ricevere campioni;
    \item un arm molto promettente riceve lo stesso numero di campioni degli altri.
\end{itemize}

Sarebbe invece desiderabile esplorare maggiormente gli arm che sembrano più promettenti.

\subsubsection{BAI è più esplorativo del MAB classico}

Nel BAI il cumulative reward ottenuto durante il testing non costituisce direttamente il criterio di valutazione.

Le slide propongono l'analogia di disporre di monete ``a perdere'': non importa perdere reward durante l'esperimento purché al termine si riesca a identificare correttamente l'arm migliore.

Anche nel BAI esiste quindi un bilanciamento tra exploration ed exploitation, ma rispetto al MAB orientato al cumulative reward la bilancia è maggiormente spostata verso l'exploration.

\subsection{Sequential Halving}

Tra le numerose strategie proposte per Fixed-Budget Best Arm Identification, le slide introducono \emph{Sequential Halving}, descritta come una strategia:
\begin{itemize}
    \item molto semplice;
    \item abbastanza efficace in diverse situazioni.
\end{itemize}

L'idea fondamentale consiste nell'alternare:
\begin{enumerate}
    \item campionamento degli arm ancora attivi;
    \item eliminazione progressiva degli arm meno promettenti.
\end{enumerate}

Il budget complessivo:
\[
T
\]
viene suddiviso in più porzioni o fasi.

La slide afferma che tali fasi hanno una certa dimensione, ma nel materiale visualizzato non viene riportata una formula esplicita per tale dimensione; pertanto non è possibile specificarla ulteriormente senza introdurre informazione esterna alle slide.

Il procedimento descritto è:
\begin{enumerate}
    \item nella prima fase, i sample vengono distribuiti uniformemente su tutti gli arm;
    \item al termine della fase si calcola, per ciascun arm, la media empirica;
    \item gli arm vengono ordinati secondo le rispettive medie empiriche;
    \item la metà degli arm che si trova nella parte inferiore del ranking viene eliminata;
    \item la fase successiva opera soltanto sugli arm superstiti;
    \item il processo continua eliminando progressivamente metà degli arm;
    \item al termine delle fasi rimane un unico arm, che viene restituito come best arm.
\end{enumerate}

Il principio è quindi quello di non sprecare progressivamente il budget sugli arm che, sulla base delle osservazioni già raccolte, appaiono chiaramente poco promettenti.

\subsection{Limiti e varianti del Best Arm Identification}

Strategie come Sequential Halving dipendono fortemente da due assunzioni:
\begin{enumerate}
    \item il numero di arm deve essere finito;
    \item il budget complessivo deve essere noto a priori.
\end{enumerate}

Le slide evidenziano diversi casi in cui tali assunzioni vengono violate.

\subsubsection{Problemi Anytime}

Tecniche di \emph{Monte Carlo Tree Search} per giochi a turni possono essere \emph{Anytime}.

In tali contesti non è necessariamente disponibile un budget prefissato prima dell'esecuzione.

Per problemi BAI Anytime, le slide osservano che viene spesso utilizzato:
\[
\mathrm{UCB1}.
\]

\subsubsection{Iperparametri continui}

Gli iperparametri da ottimizzare sono frequentemente continui.

Una discretizzazione dello spazio dei possibili valori può generare un numero molto elevato di arm, rendendo problematica l'assunzione di un piccolo insieme finito di alternative.

\subsubsection{Duelling Bandits}

Un'ulteriore variante si verifica quando un arm non produce un reward significativo se valutato isolatamente.

Il reward emerge invece dal confronto fra due arm in un \emph{match}.

Le slide indicano questo scenario con il termine:
\[
\text{\emph{duelling bandits}}.
\]

\subsubsection{Estensioni del modello classico}

Per affrontare casi che non rispettano le assunzioni del BAI classico sono stati sviluppati modelli differenti.

Le slide citano, come esempio:
\[
\mathrm{CLOP}.
\]

L'idea conclusiva è quindi che tanto il Multi-Armed Bandit classico quanto il Best Arm Identification costituiscano modelli di base dai quali derivano numerose varianti, necessarie per adattare il compromesso fra apprendimento, esplorazione e sfruttamento alle caratteristiche dello specifico problema applicativo.